\documentclass[10pt]{article}
\usepackage{icmlko}

\icmlmeta{IP 파운데이션 월드모델}{477}{영(零)은 합이 아니라 상쇄였다}{판이 낸 $\Delta \approx 0$ 을 갈랐더니 열두 도메인 중 열에서 씨앗에 안정적인 효과가 있었고, 다섯은 양수 다섯은 음수였다}

\begin{document}

\twocolumn[
\icmlheader

\begin{abstract}
직전 사이클이 이 판의 자를 의심했다. 도메인마다 주는 \emph{특기 칸}을 하나에서
둘로 늘렸을 때 판은 $+0.0008$ 로 문턱 $0.0045$ 를 못 넘었는데, 그 안에서
아이돌 도메인만 $+0.0274$ 로 크게 벌었다. 그래서 \textbf{판이 못 보는 이득이
있다}고 적었다. 그런데 그 $+0.0274$ 에는 \textbf{오차 막대가 없었다} ---
씨앗 셋의 도메인 평균을 뺀 값이었고, 그 도메인의 씨앗 SD 는 이미 $0.0236$
으로 기록돼 있었다.

이번 사이클은 팔이 아니라 \textbf{자}를 쟀다. 씨앗 열둘을 두 팔에 공유해
자 넷을 댔다 --- ① 판(유보수 가중) ② 거시판(비가중) ③ 도메인 씨앗 짝
④ 도메인 행 군집 부트. 결과는 \textbf{①②④ 판정 불능 · ③ 통과}다.
사전등록한 규칙이 그 조합에 이미 답을 정해 두었다: \emph{씨앗 잡음은 이겼고
표본 잡음은 못 이겼다.} 그래서 \textbf{``아이돌이 실제로 번다''}를
\textbf{``방향은 일관되나 49행으로는 이득이라 부를 수 없다''}로 정정한다.
거시판도 안 들인다 --- 짝 SE 가 판의 \textbf{1.78배}로 더 둔했고, 채택 조건을
측정 전에 못박아 두었다.

\textbf{그런데 진짜 결과는 도메인 표에 있었다.} 판 $\Delta$ 가 $-0.0004$ 인 그
팔에서, \textbf{열두 도메인 중 열이 $|\Delta|/\mathrm{SE} \ge 2$}였다
(최대 $4.56$). 그리고 \textbf{다섯은 양수, 다섯은 음수}다. 가중을 곱해 더하면
양수합 $+0.00174$ · 음수합 $-0.00209$ · 순합 $-0.00035$ --- 실측 판 $\Delta$
$-0.00036$ 을 재현한다. \textbf{판이 낸 영은 ``아무 일도 안 일어났다''가
아니라 ``일어난 일들이 서로를 지웠다''였다.}
\end{abstract}
]

\section{자에 오차 막대를 안 달았다}

직전 사이클의 결론은 한 문장이었다 --- \emph{제일 크게 버는 도메인이 판에서
제일 안 보인다.} 근거는 아이돌 도메인 $\Delta$ $+0.0274$ 하나였다.

그 수치에는 \textbf{잡음 바닥이 없었다.} 씨앗 셋의 도메인 평균을 뺀 값이고,
씨앗별 값도 산출물에 안 남겼다. 그런데 이 판의 문서는 이미 적고 있었다 ---
\emph{``도메인 표기값은 더 심하다 --- 씨앗 SD 가 아이돌 $0.0420$''}.
\textbf{보고한 효과가 그 도메인의 씨앗 SD 보다 작았다.}

짝(같은 씨앗)이면 그 흔들림은 상쇄된다. 문제는 \textbf{얼마나 상쇄되는지를
아무도 안 쟀다}는 것이다. 그래서 이번에는 팔이 아니라 자를 쟀다.

\claimx{바닥을 안 잰 수치 위에 다음 사이클을 쌓으면, 그 사이클은 앞
사이클의 잡음을 측정한다.}{씨앗을 늘렸을 때 $|\Delta|/\mathrm{SE}$ 가
안 변하면 이 걱정이 과했던 것이다.}

씨앗을 셋에서 열둘로 늘리자 값이 이렇게 움직였다:

\begin{center}\small
\begin{tabular}{rrrr}
\toprule
씨앗 & 짝 $\Delta$ & $|\Delta|/\mathrm{SE}$ & 양수 \\
\midrule
3 & $+0.0274$ & \textbf{1.96} & 3/3 \\
6 & $+0.0270$ & 4.00 & 6/6 \\
9 & $+0.0205$ & 3.38 & 7/9 \\
12 & $+0.0267$ & 3.70 & 10/12 \\
\bottomrule
\end{tabular}
\end{center}

$n{=}3$ 에서 \emph{양수 3/3} 은 완벽해 보이지만
$|\Delta|/\mathrm{SE}$ 는 \textbf{2 에도 못 미쳤다}. SE 를 안 냈으므로 그것을
볼 수가 없었다.

\section{자 넷}

\textbf{배선을 먼저 닫았다.} 자료 해시가 직전 사이클과 같고, $K{=}1$ 의 특기열
해시가 같으며, 재구현한 평가 경로가 정본 하네스와 \textbf{부동소수 정확히}
일치한다($0.4724867181663707$). 그리고 두 팔의 아이돌 예측이 실제로 다른지
확인했다 --- 같으면 $\Delta$ 는 통계가 아니라 항등식이고 측정이 무의미하다.

\begin{center}\small
\begin{tabular}{lrrl}
\toprule
자 & 짝 $\Delta$ & $|\Delta|/\mathrm{SE}$ & 판정 \\
\midrule
① 판(가중) & $-0.00036$ & 0.65 & 판정 불능 \\
② 거시판(비가중) & $+0.00038$ & 0.39 & 판정 불능 \\
③ 아이돌 씨앗 짝 & $+0.02673$ & \textbf{3.70} & 통과 \\
④ 아이돌 행 BCa & $+0.03458$ & $[-0.049,+0.119]$ & 판정 불능 \\
\bottomrule
\end{tabular}
\end{center}

③ 은 \emph{씨앗을 다시 뽑아도 같은가}를 묻고 ④ 는 \emph{이 51행이 다른
51행이었어도 같은가}를 묻는다. 사전등록은 \textbf{둘 다} 요구했다 --- ③ 만
이기면 씨앗 잡음만 이긴 것이고 표본 잡음은 안 본 것이기 때문이다. ④ 의 구간은
폭 $0.168$ 로 0 을 문다. 게다가 군집이 51행에 50개(병합 1)라 사실상 행
부트이고, 그것은 폭이 \textbf{과소} 방향이다 --- 관대하게 재고도 못 넘었다.

\claimx{아이돌 팔의 방향은 씨앗에 일관되지만, 49행으로는 그것을 이득이라
부를 수 없다.}{유보 행이 늘거나 군집 구조가 생기면 같은 팔이 ④ 에서도 서야
한다. 안 서면 방향 자체가 씨앗 인공물이다.}

거시판은 \textbf{안 들인다}. 채택 조건을 측정 전에 못박아 두었다 --- \emph{짝
SE 가 판 이하일 때만}. 실측은 $0.00098$ 대 $0.00055$ 로 \textbf{1.78배 더
둔했다}. 아이돌을 보이게 만드는 자는 아이돌의 잡음도 같이 보이게 만든다.
\emph{아이돌이 보이니까} 더 둔한 자를 들이는 것은 자를 결과에 맞춰 고르는
일이다.

\section{영(零)을 갈랐더니}

\begin{figure}[t]
\centering
\includegraphics[width=\columnwidth]{figs/fig1.pdf}
\caption{판 $\Delta$ 가 $-0.0004$ 인 그 팔의 내부. 왼쪽은 각 도메인이 판에
실은 몫, 오른쪽은 그 효과의 씨앗 안정성. \textbf{열둘 중 열이
$|\Delta|/\mathrm{SE}\ge2$ 이고 부호는 정확히 반씩 갈린다.}}
\label{fig:1}
\end{figure}

판이 $-0.0004$ 를 냈으므로 \emph{아무 일도 없었다}고 읽기 쉽다. 도메인별로
가르면 그렇지 않다.

\begin{center}\small
\begin{tabular}{lrrr}
\toprule
 & 짝 $\Delta$ & $|\Delta|/\mathrm{SE}$ & 양수 \\
\midrule
게임 & $+0.00515$ & \textbf{4.56} & 10/12 \\
애니 & $-0.00231$ & \textbf{4.37} & 2/12 \\
펀딩 & $-0.00393$ & 3.89 & 1/12 \\
아이돌 & $+0.02673$ & 3.70 & 10/12 \\
도서 & $-0.01502$ & 3.49 & 2/12 \\
팝업 & $-0.01412$ & 2.78 & 1/12 \\
영화 & $-0.00258$ & 2.65 & 2/12 \\
세계애니 & $+0.00307$ & 2.58 & 9/12 \\
만화 & $+0.00194$ & 2.28 & 9/12 \\
웹툰 & $+0.00372$ & 2.04 & 9/12 \\
\midrule
모바일 & $+0.00058$ & 1.11 & 8/12 \\
시장팝업 & $+0.00133$ & 0.85 & 6/12 \\
\bottomrule
\end{tabular}
\end{center}

가중을 곱해 더하면 \textbf{양수합 $+0.00174$ · 음수합 $-0.00209$ · 순합
$-0.00035$} 이고 실측 판 $\Delta$ 는 $-0.00036$ 이다. 산술이 재현된다.

\claimx{판이 낸 $\Delta \approx 0$ 은 ``효과가 없다''가 아니라 ``씨앗에
안정적인 효과 열 개가 서로를 지웠다''였다.}{같은 팔을 다시 재서 부호 무늬가
바뀌면 --- 특히 $|\Delta|/\mathrm{SE}\ge4$ 인 게임·애니의 부호가 뒤집히면 ---
이 주장은 무너진다.}

이것은 직전 사이클의 진술보다 \textbf{강한} 진술이다. 거기서는 \emph{판이 얇은
도메인을 못 본다}고 했다. 실제로는 그보다 나쁘다 --- 판은 \textbf{두꺼운
도메인의 확실한 효과들끼리도 부호가 갈리면 지운다.} 애니($-4.37\sigma$,
가중 $0.16$)와 게임($+4.56\sigma$, 가중 $0.05$)이 서로를 상쇄한다. 얇고
두껍고의 문제가 아니라 \textbf{부호가 섞인 것을 하나의 스칼라로 접는 일}
자체의 문제다.

\section{그래서 무엇이 남나}

이 사이클은 채택을 하나도 안 냈다. 남긴 것은 셋이다.

\begin{enumerate}\itemsep3pt
\item \textbf{수치 정정.} $+0.0274 \to +0.02673$(12씨앗). 값은 비슷하지만
      3씨앗일 때 $|\Delta|/\mathrm{SE}$ 가 $1.96$ 이었다는 사실이 새로 드러났다.
\item \textbf{문장 정정.} \emph{``아이돌이 실제로 번다''} $\to$
      \emph{``49행으로는 이득이라 부를 수 없다''}.
\item \textbf{새 물음.} 부호가 갈리는 도메인 효과를 하나의 가중 평균으로
      접는 것이 옳은가. 접기 전에 볼 수 있는 자 --- 예를 들어 \emph{부호 일치
      도메인 수}나 \emph{도메인별 문턱 통과 개수} --- 를 부차 자로 둘 수 있나.
      이번에 거시판은 그 후보였고 \textbf{더 둔해서 떨어졌다}. 다음 후보는
      평균이 아니라 \textbf{계수}여야 한다.
\end{enumerate}

마지막 항목이 이 실험실의 오랜 습관을 건드린다. 판 $\rho$ 하나로 채택을
결정해 온 백 노트 동안, \emph{판이 안 움직였다}고 적힌 판정들 가운데 얼마나
많은 것이 사실은 \textbf{상쇄}였는지 우리는 모른다. 그것은 다시 셀 수 있다 ---
그 실험들의 도메인별 산출물이 남아 있는 한.

\end{document}
